\section{ماژول astro\_quantizer (Astro Quantizer)}
\label{sec:module_astro_quantizer}

\subsection{معرفی و انگیزه}
این ماژول در چارچوب پژوهش حاضر پیاده‌سازی و ارزیابی شده و فاقد مقاله مستقل جداگانه در ادبیات تحقیق است.

کوانتیزه‌سازی یکنواخت وزن‌های شبکه‌های عصبی عمیق به مقادیر با طول بیت پایین (\lr{Low-Bit Quantization})، موثرترین روش جهت کاهش ردپای حافظه و افزایش بازدهی مصرف انرژی است. با این وجود، در روش‌های سنتی کوانتیزه‌سازی با طول بیت ثابت، به تمام کانال‌ها و لایه‌ها نرخ بیت یکسانی اختصاص می‌یابد؛ در حالی که در ترنسفورمرهای بزرگ، توزیع انرژی و حساسیت کانال‌های وزنی به شدت ناهمگن است و بخش اندکی از ابعاد ویژگی سهم عمده‌ای در دقت پیش‌بینی مدل ایفا می‌کنند.

موتور کوانتیزه‌سازی پویا \lr{AstroQuantizer} با هدف غلبه بر این چالش طراحی گردیده است. این چارچوب نوآورانه از سیستم تخصیص پویای نرخ بیت بر اساس اطلاعات فیشر تجربی (\lr{Fisher-Adaptive Mixed-Precision Bit Allocation}) بهره می‌برد. در این مکانیزم، کانال‌هایی که حساسیت گرادیانی یا انرژی فعال‌سازی بالاتری دارند، طول بیت بیشتری (تا سقف ۸ بیت) دریافت کرده و کانال‌های کم‌اهمیت به بیت‌های پایین‌تر (تا کف ۲ بیت) نگاشته می‌شوند، در حالی که بودجه میانگین طول بیت کل لایه در سطح هدف دقیقا ارضا می‌گردد. همچنین با استفاده از ارزیاب گرادیان مستقیم (\lr{Straight-Through Estimator - STE})، جریان گرادیان در کل شبکه بدون انسداد حفظ می‌شود.

\subsection{مبانی نظری و فرمول‌بندی ریاضی}
فرآیند کوانتیزه‌سازی در \lr{AstroQuantizer} در دو مرحله اساسی صورت می‌پذیرد:
\begin{enumerate}
    \item \textbf{تخصیص بهینه بیت بر پایه مسئله کوله‌پشتی:} فرض کنید بردار اهمیت تجربی کانال‌های خروجی با $\mathbf{s} \in \mathbb{R}^N$ مشخص شده باشد (محاسبه‌شده از طریق انرژی خروجی یا فیشر تجربی). هدف، تخصیص طول بیت $b_i \in [b_{\text{min}}, b_{\text{max}}]$ برای هر کانال $i \in \{1, \dots, N\}$ به نحوی است که مجموع حساسیت وزنی بیشینه شود:
    \begin{equation}
    \max_{\{b_i\}} \sum_{i=1}^N \mathbf{s}_i \cdot b_i \quad \text{subject to} \quad \sum_{i=1}^N b_i \le \operatorname{round}(N \cdot \bar{b}), \quad b_{\text{min}} \le b_i \le b_{\text{max}}
    \end{equation}
    که در آن $\bar{b}$ بودجه میانگین بیت هدف است. این مسئله به روش حریصانه دقیق حل می‌شود: کانال‌ها بر حسب اهمیت $\mathbf{s}_i$ به صورت نزولی مرتب شده، ابتدا به تمام کانال‌ها بیت پایه $b_{\text{min}}$ اختصاص یافته و سپس بودجه باقیمانده به کانال‌های دارای بالاترین اهمیت تا سقف $b_{\text{max}}$ تخصیص می‌یابد.
    \item \textbf{کوانتیزه‌سازی با مقیاس کانالی و ارزیاب \lr{STE}:} برای هر کانال $i$ با طول بیت تخصیص‌یافته $b_i$، محدوده گسسته $[q_{\text{min}}^{(i)}, q_{\text{max}}^{(i)}]$ و فاکتور مقیاس‌گذاری متقارن محاسبه می‌شود:
    \begin{equation}
    q_{\text{max}}^{(i)} = 2^{b_i - 1} - 1, \quad q_{\text{min}}^{(i)} = -2^{b_i - 1}, \quad \Delta_i = \frac{\max(|\mathbf{W}_{i, :}|)}{q_{\text{max}}^{(i)}}
    \end{equation}
    نگاشت کوانتیزه‌سازی و بازگشت به فضای پیوسته به صورت زیر انجام می‌پذیرد:
    \begin{equation}
    \mathbf{W}_{i, j}^{\text{quant}} = \operatorname{clamp}\left(\left\lfloor \frac{\mathbf{W}_{i, j}}{\Delta_i} \right\rceil, q_{\text{min}}^{(i)}, q_{\text{max}}^{(i)}\right) \cdot \Delta_i
    \end{equation}
    در گذر پسرو، به دلیل صفر بودن مشتق تابع پله‌ای گردکردن، مشتق معکوس توسط عملگر \lr{STE} مستقیماً عبور داده می‌شود:
    \begin{equation}
    \frac{\partial \mathcal{L}}{\partial \mathbf{W}} \approx \frac{\partial \mathcal{L}}{\partial \mathbf{W}^{\text{quant}}}
    \end{equation}
\end{enumerate}

\subsection{معماری و پیاده‌سازی}
این ماژول در فایل \lr{neuralops/compression\_engine/astro\_quantizer.py} قرار دارد و در سیستم با شناسه‌های \lr{astro} و \lr{astro\_quantizer} ثبت گردیده است:
\begin{itemize}
    \item \textbf{کلاس تابع اتوگرد \lr{STEFakeQuantize}:} تابعی ارث‌بری‌شده از \lr{torch.autograd.Function} که عملیات نگاشت گسسته و بازگشت را با پشتیبانی از بردار بیت‌های ناهمگن پیاده کرده و مشتق را در مسیر پسرو بدون تغییر منتقل می‌کند.
    \item \textbf{کلاس \lr{AstroLinear}:} لایه خطی جایگزین که وزن کوانتیزه‌شده شبیه‌سازی‌شده (\lr{fake-quantized}) را به عنوان بافر یا پارامتر ذخیره کرده و سازگاری کامل با لایه‌های \lr{Conv1D} در مدل‌های زبانی مانند \lr{GPT-2} را تامین می‌نماید.
    \item \textbf{کلاس موتور \lr{AstroQuantizer}:} هماهنگ‌کننده کالیبراسیون و اجرای الگوریتم کوله‌پشتی برای توزیع بیت روی ماتریس وزن‌ها با رعایت فیلترهای استثنای لایه‌های حساس.
\end{itemize}

\subsection{نتایج تجربی}
نتایج ارزیابی آزمایشگاهی \lr{astro\_quantizer} بر روی مدل‌های مرجع در جدول \ref{tab:benchmark_astro_quantizer} به تفکیک ارائه شده است:

\begin{table}[H]
\centering
\small
\caption{نتایج بنچ‌مارک سخت‌افزاری و زبانی فشرده‌ساز \lr{astro\_quantizer} روی مدل‌های پایه}
\label{tab:benchmark_astro_quantizer}
\begin{tabular}{lccccc}
\toprule
\textbf{مدل پایه \lr{(Base Model)}} & \textbf{دقت \lr{(Accuracy)}} & \textbf{شاخص \lr{F1}} & \textbf{حافظه \lr{(Peak VRAM)}} & \textbf{توان \lr{(Throughput)}} & \textbf{وضعیت} \\
\midrule
\lr{DistilBERT (SST-2)} & \lr{90.48\%} & \lr{0.9064} & \lr{368.29 MB} & \lr{171,343.82 tok/s} & \lr{PASS} \\
\lr{DistilBERT (MNLI)} & \lr{81.80\%} & \lr{0.8174} & \lr{457.35 MB} & \lr{184,967.17 tok/s} & \lr{PASS} \\
\lr{RoBERTa (SST-2)} & \lr{94.04\%} & \lr{0.9421} & \lr{514.29 MB} & \lr{12,245.45 tok/s} & \lr{PASS} \\
\lr{GPT-2 (WikiText-2)} & \lr{0.00\%} & \lr{0.0000} & \lr{0.00 MB} & \lr{0.00 tok/s} & \lr{CRASH} \\
\lr{Qwen} & -- & -- & -- & -- & \lr{NOT\_RUN} \\
\bottomrule
\end{tabular}
\end{table}

همان‌طور که در جدول مشاهده می‌شود، ماژول \lr{astro\_quantizer} در مدل‌های انکودری \lr{RoBERTa} و \lr{DistilBERT} بالاترین دقت‌های کل موتور فشرده‌سازی را به ثبت رسانده است: در مدل \lr{RoBERTa} دقت قابل توجه \lr{94.04\%} با شاخص F1 معادل \lr{0.9421} حاصل شده و در \lr{DistilBERT} روی دادگان SST-2 دقت \lr{90.48\%} و روی MNLI دقت \lr{81.80\%} به دست آمده است. این ارقام نشان‌دهنده اثربخشی خارق‌العاده تخصیص غیریکنواخت بیت بر پایه اطلاعات فیشر است. روی مدل‌های خانواده \lr{Qwen} تستی صورت نگرفته و در وضعیت \lr{NOT\_RUN} قرار دارد.

\subsection{محدودیت‌ها و تحلیل شکست}
تحلیل رخدادهای بنچ‌مارک و بررسی کد منبع موارد زیر را نمایان می‌سازد:
\begin{itemize}
    \item \textbf{رخداد خطای زمان اجرا در بنچ‌مارک دیکودری \lr{GPT-2}:} در ارزیابی مدل \lr{GPT-2} بر روی مجموعه دادگان WikiText-2، بنچ‌مارک دچار وضعیت کرش گردید (\lr{CRASH}). بررسی ردپای خطا حاکی از بروز استثنای متغیر محلی در مرحله ارزیابی دقت است (\lr{UnboundLocalError: cannot access local variable 'majority\_count'}). این باگ نرم‌افزاری ناشی از شرط‌بندی نادرست بر روی ابعاد لاجیت‌های توکن‌های دنباله دیکودری در اسکریپت درایور بنچ‌مارک بوده و ناشی از ماهیت ریاضی خود کوانتیزه‌ساز نیست.
    \item \textbf{پیچیدگی شتاب‌بخشی سخت‌افزاری در پردازش بیت‌های ناهمگن:} اگرچه کوانتیزه‌سازی شبیه‌سازی‌شده سربار حافظه تئوریک را کاهش می‌دهد، اما اجرای عملیاتی کانال‌های با طول بیت‌های متفاوت (مثلاً ۲ بیت در کنار ۸ بیت) روی کرنل‌های \lr{GEMM} یکپارچه پردازنده‌های گرافیکی نیازمند کتابخانه‌های تخصصی چیدمان بیت (\lr{Bit-Packing}) است.
\end{itemize}
